%%% RnE 보고서 관련 요약문 등 %%%

%%% 연구결과 요약문 및 과제 수행 과정의 표에 들어가는 내용을 쓰는 곳이다.

%%%--------------------------------------
\newcommand{\summaryGoal}{%%% 연구 필요성 및 목적 쓰는 곳
\jiwon[1] %% todo: 여기의 \jiwon[1]를 삭제하고 "연구의 필요성 및 목적"을 쓸 것 
% "연구 필요성 및 목적 쓰는 곳 끝.
}
%%%---------------------------------------


\newcommand{\summaryResult}{%%% 연구내용 및 연구결과 쓰는 곳
\jiwon[2-3] %% todo: 여기의 \jiwon[2-3]을 삭제하고 "연구내용 및 연구결과"를 쓸 것 
% 연구내용 및 연구결과 쓰는 곳 끝.
}

%%%---------------------------------------
\newcommand{\summaryEffect}{%%% 연구 기대효과 쓰는 곳
\jiwon[4] %% todo: 여기의 \jiwon[4]를 삭제하고 "연구 기대효과"를 쓸 것 
% 연구 기대효과 쓰는 곳 끝.
}

%%%---------------------------------------


\newcommand{\summaryExpert}{%%% 전문가 피드백 활용 과정 쓰는 곳
% 전문가 피드백 활용을 받지 않았다면, 아래 줄의 \jiwon[5]를 삭제하고 "해 당 없 음"이라고 쓴다.
\jiwon[5] %% todo: 전문가 피드백 활용과정에 쓸 내용이 있으면, \jiwon[5]를 지우고 내용을 쓰면 된다.
% 전문가 피드백 활용 과정 쓰는 곳 끝.
}

%%%---------------------------------------


\newcommand{\summaryProblemSolving}{%%% 문제해결 과정 쓰는 곳
\jiwon[6-7] %% todo: 여기의 \jiwon[6-7]을 삭제하고 "문제해결 과정"을 쓸 것 
% 문제해결 과정 쓰는 곳 끝.
}

%%%---------------------------------------
\newcommand{\summaryFirstMember}{%%% 팀원1 내용 쓰는 곳
\jiwon[8] %% todo: 여기의 \jiwon[8]을 삭제하고 "팀원1 "의 내용을 쓸 것 
% 팀원1 내용 쓰는 곳 끝.
}

%%%---------------------------------------

% \newcommand{\summarySecondMember}{%%% 팀원2 내용 쓰는 곳


% % 팀원2 내용 쓰는 곳 끝.
% }


%%%---------------------------------------
\newcommand{\summaryLearning}{%%% 과제 수행 과정에서의 학습 쓰는 곳
\jiwon[9] %% todo: 여기의 \jiwon[9]를 삭제하고 "과제 수행 과정에서의 학습"의 내용을 쓸 것 
% 과제 수행 과정에서의 학습 쓰는 곳
}